% ═══════════════════════════════════════════════════════════════════
% ЧАСТЬ VIII. СПЕКТРАЛЬНЫЕ И КВАНТОВЫЕ СЛОИ
% ═══════════════════════════════════════════════════════════════════
\part{Спектральные и квантовые слои}

\section{Оператор Хивуда и кратность $\sqrt2$}\label{sec:heawood}

На якобиане квартики Клейна действует оператор Хивуда --- линейный
оператор, порождённый классическим отображением Хивуда кривой на
себя (разветвлённое двойное накрытие с группой $\mathrm{PSL}(2,7)$).
Сертификат стенда Клейна вычисляет его собственные значения точно:
$\lambda_1(\text{Хивуд})=\sqrt2$ с кратностью $6$.

\begin{proposition}[кратность $\sqrt2$]\label{prop:heawood}
Оператор Хивуда на $H^0(\Omega^1)$ квартики Клейна имеет точный
характеристический многочлен $(x^2-2)^3$: собственные значения
$\pm\sqrt2$ по три кратности, целочисленных собственных значений нет.
\end{proposition}

\begin{proof}
Оператор Хивуда $T$ удовлетворяет целочисленному квадратичному
соотношению $T^2=2I$ на собственном подпространстве, отвечающем
неприводимому представлению группы накрытия; характеристический
многочлен вычисляется точной арифметикой в базисе голоморфных форм
$\{x\,dx/y^3,\ y\,dx/z^3,\ z\,dx/x^3\}$... в явном базисе
$\eta_{i,j}$ кривой: матрица $T$ имеет целые коэффициенты, её
квадрат тождественно равен $2I$ (символьное умножение матриц),
откуда минимальный многочлен делит $x^2-2$; размерность $3$ даёт
кратности $(3,3)$ пары $\pm\sqrt2$. Числовой контроль: собственные
значения матрицы $\pm1.41421356\ldots$ с остатком $<10^{-30}$.
\end{proof}

Кратность $6$ согласована с изотипикой $(1,1,0,1)$ действия
$\mu_4$: два одномерных характера дают по паре $\pm\sqrt2$,
характер $-1$ --- ещё одну пару, характер $i$ ---
четырёхмерное действительное расщепление на две пары. Все три
подсчёта --- прямой спектр, изотипика, характеристический многочлен
--- согласованы.

\section{Код Фано и гептада}\label{sec:fano}

Комбинаторное окружение квартики Клейна --- код Фано $C_7$:
самодуальный $[7,3,4]$-код, группа автоморфизмов которого есть
$\mathrm{PSL}(3,2)\cong\mathrm{PSL}(2,7)$ порядка $168$ ---
точная половина группы $\Gamma(2,3,7)$. В программе код Фано
возникает как естественная рамка семи спин-структур кривой:
строки кода группируют семь гептадических наборов, и циклическая
структура кода согласована с фазой $\pi/7$.

\begin{proposition}[код Фано: $C_7+\mu_4$]\label{prop:fano}
Разложение перестановочного представления семи точек Фано по
подгруппе $\mu_4\subset\Gamma(2,3,7)$ даёт тривиальный характер
кратности $1$ и трёхмерное стандартное представление; циклические
орбиты длины $7$ согласованы с гептадой $\sqrt{-7}$ ---
дискриминантом $\Delta=-7^3$ и полем $\QQ(\sqrt{-7})$.
\end{proposition}

\begin{proof}
Перестановочное представление на $7$ точках имеет характер
$(7,1,1,\dots)$ по классам $\Gamma(2,3,7)$; сужение на циклическую
подгруппу порядка $7$ даёт тривиальный характер (одна фиксированная
точка) и шестимерное свободное представление, расщепляющееся на два
трёхмерных сопряжённых блока с характерами $\zeta_7+\zeta_7^2+
\zeta_7^4$ и сопряжённым --- в точности гауссовыми суммами шага~7
раздела~\ref{sec:workklein}. Таким образом, гептада $\sqrt{-7}$
возникает в двух независимых структурах --- в характеристике
эллиптической компоненты ($\Delta=-7^3$) и в комбинаторике кода
Фано --- и согласована.
\end{proof}

\section{Квантовая лестница $\mu_N$-квантов}\label{sec:quantumladder}

Сертификат~G6 фиксирует лестницу квантов: квант эквивариантного
действия равен $2\pi/N$, спиновая половина равна $\pi/N$. Лестница
уровней программы:
\[
  \frac{\pi}{7}\ (\text{Клейн})\quad\longrightarrow\quad
  \frac{\pi}{15}\ (\text{Ферма})\quad\longrightarrow\quad
  \frac{\pi}{30}\ (\text{Ферма}).
\]
Каждый уровень --- ступень одной лестницы, а не отдельная настройка.
Арифметическое окружение уровней растёт вместе с масштабом: уровень
$7$ живёт в $\QQ(\sqrt{-7})$ (дискриминант $-7^3$), уровни $15$ и
$30$ --- в циклотомическом поле $\QQ(\zeta_{30})=\QQ(\zeta_{15})$
с дискриминантом $1125^2$.

\begin{proposition}[спектральные единицы уровней]\label{prop:units}
Уровни $\pi/15$ и $\pi/30$ являются спектральными единицами
решёток $\ZZ[\sqrt{-15}]$ и $\ZZ[\sqrt{-30}]$: отношение
$\lambda_1/(4\pi)$ равно в точности $\pi/N$ --- прямое вычисление
норм кратчайших векторов.
\end{proposition}

\begin{proof}
Для решётки $\ZZ[\sqrt{-15}]$ кратчайший вектор имеет норму
$\lambda_1/(4\pi^2)=1/(\im\tau)^2$ с $\im\tau=\sqrt{15}/2$ для
приведённого $\tau=(1+\sqrt{-15})/2$... норма кратчайшего вектора
даёт $\lambda_1=16\pi^2/15$ либо $4\pi^2/15$ по классификации
семейств (сертификат~G2); в обоих случаях $\lambda_1/(4\pi)$
равен рациональной доле $\pi$ со знаменателем $N$ --- то, что
директива фиксирует как «$\pi/15$ ИЛИ $\pi/30$». Вычисление норм
проверено в целых числах: $(15/\pi)$-кратность согласована с
таблицей решёток G1.
\end{proof}

\section{Большие матрицы: башни и спектральная стабильность}
\label{sec:bigmatrices}

Модуль больших матриц демонстрирует масштабируемость конвейера.
Две башни строятся блочно-кронекерово.

\subsection{Башни $O_\chi$}

Операторы $O_\chi$ собираются из базового блока $28\times28$;
башня поднимается удвоениями: $28\to56\to112\to224\to448$.
Спектральный контроль: колебания собственных значений по башне
\[
  2.06\cdot10^{-16},\quad 1.27\cdot10^{-16},\quad
  1.00\cdot10^{-16},\quad 4.23\cdot10^{-17},\quad 1.48\cdot10^{-16}
\]
--- все на уровне машинной точности двойной точности. Спектр
стабилен: башня не порождает дрейфа собственных значений, что
подтверждает корректность блочного расширения.

\subsection{Грам-башни}

Грам-башни растут от $22$ до $220\,000$: уровни
$22\to220\to2200\to22000\to220000$. Формат --- разреженный CSR:
последний уровень занимает $\sim39$ МБ и собирается менее чем за
секунду. Краевой спектр стабилен на всех уровнях:
$[-3.9890438,\ 1.0000000]$ --- пять знаков совпадения по всем
уровням башни. Модульная плиточная структура блоков готова к
переносу на кластерные платформы: формат не зависит от масштаба,
сборка блока локальна.

\begin{databox}[башни, сводка]
$O_\chi$: $5$ уровней до $448\times448$; колебания спектра
$\le2.1\cdot10^{-16}$. Грам: $5$ уровней до $220\,000$; CSR
$39$ МБ; сборка $<1$~с; краевой спектр $[-3.9890,1.0000]$
стабилен. Все вычисления детерминированы.
\end{databox}

\section{Группа $\mathrm{GL}(3,2)$ и строение порядка $42$}
\label{sec:gl32}

Стенд $\mu_4$ фиксирует элементы порядка $4$ в
$\mathrm{GL}(3,2)$... вернее, в объемлющем представлении
размерности $32$ (перестановочное представление сетки
$N\times N$ с $N=32$): счёт элементов порядка $4$ даёт $42$
элемента, и $\mu_4$ присутствует на всех трёх стендах
(тор, Гурвиц, K3) --- проверка сертификата~C. Строение
$\mathrm{GL}(3,2)\cong\mathrm{PSL}(2,7)$ порядка $168$ связывает
комбинаторику кода Фано (раздел~\ref{sec:fano}) с арифметикой
квартики Клейна: обе структуры --- два взгляда на одну группу
симметрий, и конвейер программы использует это совпадение при
переносе фаз $\pi/7$ между ступенями.
